\documentclass[11pt,a4paper]{article}
\usepackage[a4paper,margin=20mm,headheight=15pt,headsep=8mm,footskip=10mm]{geometry}
\usepackage{fontspec}
\setmainfont{Liberation Serif}
\setsansfont{Liberation Sans}
\usepackage{polyglossia}
\setdefaultlanguage{russian}
\usepackage{amsmath,amssymb,mathtools}
\usepackage{graphicx}
\usepackage{xcolor}
\usepackage{microtype}
\usepackage{fancyhdr}
\usepackage[most]{tcolorbox}
\usepackage{enumitem}
\usepackage{titlesec}
\usepackage{caption}
\usepackage{hyperref}
\usepackage{needspace}
\definecolor{navy}{HTML}{18344F}
\definecolor{teal}{HTML}{287D82}
\definecolor{soft}{HTML}{F2F6F8}
\definecolor{muted}{HTML}{566575}
\hypersetup{colorlinks=true,linkcolor=navy,urlcolor=teal}
\setlength{\parindent}{0pt}
\setlength{\parskip}{4pt plus 1pt}
\setlist{nosep,leftmargin=1.5em}
\titleformat{\section}{\sffamily\Large\bfseries\color{navy}}{}{0pt}{}[\vspace{1pt}\color{teal}\titlerule]
\titleformat{\subsection}{\sffamily\bfseries\color{navy}}{}{0pt}{}
\titlespacing*{\section}{0pt}{10pt}{6pt}
\titlespacing*{\subsection}{0pt}{7pt}{3pt}
\pagestyle{fancy}\fancyhf{}
\fancyhead[L]{\small\sffamily УМФ \textbullet\ Задание 1}
\fancyhead[R]{\small\sffamily Волна от сферической полости}
\fancyfoot[C]{\small\color{muted}\thepage}
\renewcommand{\headrulewidth}{0.4pt}\renewcommand{\headrule}{\hbox to\headwidth{\color{teal}\leaders\hrule height \headrulewidth\hfill}}
\newtcolorbox{answerbox}{enhanced,breakable,colback=soft,colframe=teal!70!black,boxrule=0.8pt,arc=1.5mm,left=3mm,right=3mm,top=2mm,bottom=2mm,title=\textsf{Ответ},fonttitle=\bfseries\color{navy}}
\newtcolorbox{notebox}{enhanced,breakable,colback=gray!5,colframe=gray!45,boxrule=0.5pt,arc=1mm,left=2.5mm,right=2.5mm,top=1.5mm,bottom=1.5mm}
\captionsetup{font=small,labelfont={bf,color=navy},skip=5pt}
\begin{document}

\begin{center}
{\sffamily\small\color{teal}\bfseries РАСЧЁТНАЯ РАБОТА}\par\vspace{3pt}
{\sffamily\LARGE\bfseries\color{navy} Волна сжатия от сферической полости}\par\vspace{5pt}
{\small Численное решение МКР и аналитическое решение через задачу Коши для ОДУ}
\end{center}
\vspace{2mm}
\begin{notebox}
\textbf{Постановка.} После скачкообразного приложения внутреннего давления требуется найти радиальное перемещение $u(r,t)$ в однородной изотропной линейно-упругой среде двумя способами: методом конечных разностей второго порядка и аналитически. Используем обозначения из лекций: баланс импульса, закон Гука и уравнение Ламе.
\end{notebox}

\section{Уравнение движения и исходные данные}

Для среды без объёмных сил и малых деформаций лекционные соотношения имеют вид
\begin{equation*}
\rho\,\mathbf u_{tt}=\operatorname{div}\boldsymbol\sigma,\qquad
\boldsymbol\varepsilon=\tfrac12\bigl(\nabla\mathbf u+(\nabla\mathbf u)^T\bigr),\qquad
\boldsymbol\sigma=\lambda\,\operatorname{tr}(\boldsymbol\varepsilon)\mathbf I+2\mu\boldsymbol\varepsilon.
\end{equation*}
После подстановки закона Гука в баланс импульса получаем уравнение Навье--Ламе:
\begin{equation*}
\rho\mathbf u_{tt}=(\lambda+\mu)\nabla(\operatorname{div}\mathbf u)+\mu\Delta\mathbf u.
\end{equation*}
При сферической симметрии $\mathbf u=u(r,t)\mathbf e_r$ и $\operatorname{div}\mathbf u=u_r+2u/r$. Поэтому радиальное уравнение движения принимает вид
\begin{equation}\label{eq:wave}
 u_{tt}=C_p^2\left(u_{rr}+\frac{2}{r}u_r-\frac{2}{r^2}u\right),\quad r\ge R,\qquad C_p=\sqrt{\frac{\lambda+2\mu}{\rho}}.
\end{equation}

Заданы $E=5.4\cdot10^9$ Па, $\mu=2.0\cdot10^9$ Па, $\rho=1800$ кг/м$^3$, $R=1$ м, $P_0=10^8$ Па, $h=0.01R=0.01$ м и $t_*=0.01$ с. Давление $p(t)=P_0H(t)$, перемещение положительно наружу.

Из $E=2\mu(1+\nu)$ находим $\nu=0.35$, затем
\[
\lambda=\frac{2\mu\nu}{1-2\nu}=4.6667\cdot10^9\ \text{Па},\qquad
C_p=\sqrt{\frac{\lambda+2\mu}{\rho}}=2194.2686\ \text{м/с}.
\]

Начальные условия соответствуют покоящейся среде: $u(r,0)=u_t(r,0)=0$. Для радиальных напряжений закон Гука даёт
\[
\sigma_{rr}=(\lambda+2\mu)u_r+\frac{2\lambda}{r}u.
\]
С учётом того, что сжатие соответствует отрицательному радиальному напряжению, условие на поверхности полости имеет вид
\begin{equation}\label{eq:bc}
(\lambda+2\mu)u_r(R,t)+\frac{2\lambda}{R}u(R,t)=-P_0H(t).
\end{equation}

\section{Метод конечных разностей}

Вводим равномерную сетку $r_i=R+ih$, $t_n=n\tau$ и выбираем заданный в работе шаг времени $\tau=h/C_p$. Число Куранта равно $C_p\tau/h=1$. Центральные разности второго порядка приводят к явной трёхслойной схеме
\begin{equation}\label{eq:scheme}
u_{i}^{n+1}=2u_i^n-u_i^{n-1}+C_p^2\tau^2 L_i^n,
\end{equation}
где
\[
L_i^n=\frac{u_{i+1}^n-2u_i^n+u_{i-1}^n}{h^2}
+\frac{u_{i+1}^n-u_{i-1}^n}{r_i h}-\frac{2u_i^n}{r_i^2}.
\]
Первый член аппроксимирует $u_{rr}$, второй соответствует $2u_r/r$, третий -- $-2u/r^2$. Все производные в гладкой части области аппроксимируются со вторым порядком по пространству и времени.

\subsection{Граничный и начальный слои}
Производную в условии \eqref{eq:bc} аппроксимируем центральной разностью. Введя фиктивный узел $u_{-1}^n$, получаем формулу границы
\begin{equation}\label{eq:ghost}
u_{-1}^n=u_1^n+\frac{2h}{\lambda+2\mu}\left(P_0+\frac{2\lambda}{R}u_0^n\right),\qquad t_n\ge0.
\end{equation}
Так как на начальном слое $u_i^0=0$ и $u_t(r_i,0)=0$, первый временной слой рассчитывается по разложению Тейлора:
\[
u_i^1=u_i^0+\tau u_t(r_i,0)+\frac{\tau^2C_p^2}{2}L_i^0.
\]
Внешняя граница расчётной области выбрана за фронтом волны, чтобы не влиять на результат к моменту $t_*$.

\section{Аналитическое решение через ОДУ}

Для исходящей сферической волны вводим задержанное время
\[
s=t-\frac{r-R}{C_p}.
\]
Решение представим в форме
\begin{equation}\label{eq:analytic_u}
u(r,t)=\frac{B'(s)}{C_p r}+\frac{B(s)}{r^2},\qquad u(r,t)=0\quad\text{при }s<0.
\end{equation}
Подстановка в граничное условие \eqref{eq:bc} приводит к задаче Коши
\begin{equation}\label{eq:ode}
B''+aB'+bB=\frac{P_0R}{\rho},\quad B(0)=B'(0)=0,
\qquad a=\frac{4\mu}{\rho C_pR},\quad b=\frac{4\mu}{\rho R^2}.
\end{equation}
Обозначим $\omega=\sqrt{b-a^2/4}$ и $B_\infty=P_0R^3/(4\mu)$. Решение ОДУ при $t\ge0$:
\begin{align*}
B(t)&=B_\infty\left[1-e^{-at/2}\left(\cos\omega t+\frac{a}{2\omega}\sin\omega t\right)\right],\\
B'(t)&=\frac{P_0R}{\rho\omega}e^{-at/2}\sin\omega t.
\end{align*}
Вместе с формулой \eqref{eq:analytic_u} это даёт аналитическое перемещение для сравнения с МКР.

\Needspace{16\baselineskip}
\section{Результаты}

Для заданной сетки $\tau=h/C_p=4.557327\cdot10^{-6}$ с. Момент $t_*=0.01$ с не совпадает с целым числом шагов, поэтому значение численного решения в этот момент получено линейной интерполяцией между соседними временными слоями. Расчётная область содержит $N=2295$ интервалов, внешняя граница находится при $r_N=23.95$ м, а фронт волны -- при $r_f=R+C_pt_*=22.94269$ м.

\begin{center}
\includegraphics[width=\textwidth]{task1_displacement.png}
\captionof{figure}{Радиальное перемещение при $t=0.01$ с: аналитическое решение и МКР; справа показана область фронта волны.}
\end{center}

Максимальная абсолютная разность двух решений на основной сетке равна $5.65\cdot10^{-6}$ м, что составляет около $0.0452\%$ от максимального аналитического перемещения. Дополнительная проверка при последовательном уменьшении шага вдвое дала $E_\infty=6.51\cdot10^{-6}$, $1.63\cdot10^{-6}$ и $4.07\cdot10^{-7}$ м; ошибка уменьшается примерно в четыре раза, что согласуется со вторым порядком схемы.

\begin{answerbox}
\textbf{Искомый профиль построен обоими способами.} На графике решения практически совпадают; максимальное расхождение на заданной сетке составляет $5.65\cdot10^{-6}$ м (около $0.0452\%$). Проверка на сгущающихся сетках подтверждает второй порядок точности МКР.
\end{answerbox}

\end{document}